% Table source-order #05: Absence of repeated extrema among image-free modes. \textbf{Absence of
\begin{table}[tp]
\centering
\caption{Absence of repeated extrema among image-free modes. Absence of repeated extrema: read the name literally. A mode `reaches the bar' on a metric when it is the extreme (highest or lowest) in ≥7 of 8 cells (87.5\%). Over 26 metrics the four image-free modes reach it on none. \textbf{Caution:} This is not a separability test. \texttt{rank\_consistency()} only counts which mode attains each metric's max/min; a mode that is consistently \emph{second} (and strongly distinguishable) never reaches the bar. Establishing non-separability would need pairwise equivalence margins with task-clustered intervals, which this does not do. \textbf{Caution:} The threshold got stricter when cells grew, and an earlier caption said otherwise: ≥7/8 is 87.5\%, not the 83\% it claimed, and the six-cell ≥5/6 it says it matches is 83.3\%: a +4.2pp shift, chosen after the cell count changed. Carrying the literal numerator (≥5/8 = 62.5\%) instead would let DOM+stext clear two metrics and this negative would appear to break, so the choice is load-bearing and is disclosed rather than defended. Source: \texttt{per\_mode\_four\_dimension\_profile\_with\_wa.json}}
\label{tab:t05}
\footnotesize
\setlength{\tabcolsep}{3pt}
\begin{tabular}{@{}ll@{}}
\toprule
mode & metrics reaching the bar \\
\midrule
Vision & 9 \\
SoM & 5 \\
DOM & 0 \\
DOM+\allowbreak{}stext & 0 \\
DOM+\allowbreak{}sprompt & 0 \\
SoM-image & 0 \\
\bottomrule
\end{tabular}
\end{table}
